Suponha que \(2n\) casas de um tabuleiro \(n \times n\) estão marcadas. Mostre que para algum \(k > 1\) é possível selecionar \(2k\) casas marcadas distintas, digamos \(a_1, \dots, a_{2k}\), de tal forma que \(a_{2i-1}\) e \(a_{2i}\) estão na mesma linha, \(a_{2i}\) e \(a_{2i+1}\) estão na mesma coluna, para todo \(i\), com os índices tomados módulo \(2k\).